% r03_fill_K2b: proof of (K_0), part III: the equal-spacing propagator package at zero gauge field.

\subsubsection*{The equal-spacing propagator package}

Unit scale of level $r\ge2$, infinite lattice, $A=0$; $V_\eta=R^{-1}[-\frac\pi\eta,\frac\pi\eta]^3$, $V_1=R^{-1}[-\pi,\pi]^3$; $u_{r-1}$ is the common symbol of the $(r-1)$-step rotated centred average; $u_X(q)=\frac1{125}\sum_{x\in B_0}\cos(\eta'q\cdot x)$ is the tile symbol and $u_Y(q)=\prod_\alpha W(\eta'q\cdot f_\alpha)$, $W(t)=\frac15\sum_{j=-2}^2\cos jt$, the cube symbol; $\Delta'_X(q)=4\eta'^{-2}\sum_\mu\sin^2\frac{\eta'q_\mu}2+m_r^2$ and $\Delta'_Y$ its rotated analogue, $m_r=m\delta_r$ the running mass; $L^F_2$ are the $125$ alias classes; $A^F$ is the coarse multiplier of flow $F$ (the symbol of $\Delta^{F(r)}$), $A^F\le a$; $\Delta_R$ is the rotated Laplacian at spacing $\eta$; $J_s\defeq\int_{V_\eta}|u_{r-1}(q)||q|^{-s}dq$, bounded uniformly in $r$ for $s=1$. Constants depend only on $a$ and are uniform in $r$, $k$, the torus and the mass $m_r^2\ge0$.

\begin{proposition}[free core at equal spacing]\label{rot:prop:fc}
$A^F(p)^{-1}=a_{r-1}^{-1}+\sum_{l_1\in\Gamma_r^*/\Gamma_1^*}u_{r-1}(p+l_1)^2[\eta^2a^{-1}+S_F(p+l_1)]$ with $S_F(q)\defeq\sum_{l\in L^F_2}u_F(q+l)^2/\Delta'_F(q+l)$, and $\sup_{V_\eta}|S_X-S_Y|\le0.198\,\eta^2$. Hence for all $r\ge2$, all $p$ and all $m_r\ge0$,
\[
|A^X(p)-A^Y(p)|\le0.20\,\eta^2A^X(p)A^Y(p)\le0.20\,a\,\eta^2A^Y(p),
\]
and the same for the photon (three copies at mass $\lmu\delta_r$).
\end{proposition}
\begin{proof}
The formula is the Woodbury identity for a composite average whose symbol is $u_{r-1}u_F$; the terms $\eta^2/a$ cancel in the difference and $\sum_{l_1}u_{r-1}^2=1$ gives the sup bound. Scale to $\eta=1$, $\eta'=\frac15$. \emph{Zones:} on $V_1$, $\Delta'_X\ge0.93718|q|^2+m^2$ and $\Delta'_Y\ge0.96753|q|^2+m^2$ (monotonicity of $\sin x/x$), while $\Delta'_F\le|q|^2+m^2$ and $(|q|^2+m^2)-\Delta'_F\le\eta'^2|q|^4/12$. \emph{Second moments:} $1-u_X\le\frac{\eta'^2}2q^{\mathrm T}M_0q$ with $M_0=\diag(\frac{52}{25},\frac{52}{25},2)$ (exact enumeration of the tile), so $1-u_X^2\le\frac{52}{25}\eta'^2|q|^2$; $1-W(t)\le t^2$ gives $1-u_Y^2\le2\eta'^2|q|^2$. The main terms ($l=0$) then differ by at most $0.0924$. \emph{Aliases:} every non-zero alias $q+l$ has a frame component of modulus in $[\pi,5\pi]$, where $100\sum\sin^2(r_\mu/10)\ge9.549$ (concavity of $t\mapsto\sin^2(\sqrt t/10)$ on $[0,25\pi^2]$); with Parseval $\sum_lu_F(q+l)^2=1$, $0\le\mathrm{Al}_F\le(1-u_F^2)/9.549\le0.1047$. Total $0.198$. Finally $A^F\le a_{r-1}\le a$. (The true constant is about $0.005$: both first moments vanish by inversion symmetry and the second moments differ by $4\%$.)
\end{proof}

\begin{lemma}[external lines]\label{rot:lem:H}
For $x'\in x+\eta'B_0$, $x''=\tau x'\in Q_Y(x)$ ($x\in\Gamma_1$) and $z\in\Gamma_r$,
\[
H^F(x^F,z)=(2\pi)^{-3}\int_{V_\eta}dq\,e^{\mathrm iq\cdot(x-z)}u_{r-1}(q)M^F(q)A^F(q),\qquad M^F(q)\defeq\sum_{l\in L^F_2}\frac{e^{\mathrm i(q+l)\cdot r_F}u_F(q+l)}{\Delta'_F(q+l)},
\]
with $r_X=x'-x$, $r_Y=x''-x$, $x^X=x'$, $x^Y=x''$, and
\[
|H^X(x',z)-H^Y(x'',z)|\le C_EJ_1\,\eta,\qquad |H^X(x',z)-H^Y(x'',z)|\le C(C_EJ_1\eta)^{1/2}e^{-\delta|x-z|/2},
\]
$C_E=345a+0.2a(C_*/20+6.9a)$, $C_*=\max\{43.6(1+\sqrt3\pi)^2,\frac{\pi^2}4a(1+\frac1\pi)^2\}$.
\end{lemma}
\begin{proof}
The representation: $H=D^{-1}Q^*\Delta^{(r)}$ (Woodbury), $Q^*$ has symbol $u_Fu_{r-1}$; decompose the fine Brillouin zone into the translates $V_\eta+l$, $l\in L^F_2$, and use the $\Gamma_1^*$-periodicity of $u_{r-1}$, $A^F$ and $e^{-\mathrm iq\cdot z}$, with $e^{\mathrm il\cdot x}=1$. \emph{Bounds on $V_\eta$}: for both flows $|M^F-\Delta_R^{-1}|\le172.1\,\eta|q|^{-1}$. For $Y$: $|e^{\mathrm iq\cdot s}u_Y(q)-1|\le|q||s|+(1-u_Y)\le0.910\,\eta|q|$ ($|s|\le2\sqrt3\eta'$, $1-u_Y\le\eta'^2|q|^2$), divided by $\Delta'_Y\ge\frac4{\pi^2}|q|^2$ gives $2.25\,\eta|q|^{-1}$; $|\Delta'^{-1}_Y-\Delta_R^{-1}|\le\min\{0.51\eta^2,\frac{\pi^2}4|q|^{-2}\}\le1.12\,\eta|q|^{-1}$; the $124$ aliases have $\Delta'_Y(q+l)\ge4/\eta^2$ and $|u_Y|\le1$, so their sum is $\le31\eta^2\le168.7\,\eta|q|^{-1}$ on $V_\eta$. For $X$ the same three bounds are $2.37+0.73+169$. Hence $|M^X-M^Y|\le345\,\eta|q|^{-1}$. \emph{Multipliers}: $E\defeq M^XA^X-M^YA^Y=(M^X-M^Y)A^X+M^Y(A^X-A^Y)$; the first term is $\le345a\eta|q|^{-1}$; for the second use \Cref{rot:prop:fc}, $|M^Y|\le\Delta_R^{-1}+172.1\eta|q|^{-1}$ and $\Delta_R^{-1}A^Y\le C_*(1+|q|)^{-2}$ (on $V_1$ from $A^Y(q)^{-1}\ge u_{r-1}^2u_Y^2\Delta'^{-1}_Y$ with $u_{r-1}^2\ge(2/\pi)^6$, $u_Y\ge0.95$, $\Delta'_Y/\Delta_R\le\pi^2/4$; off $V_1$ from $|q|\ge\pi$), so $|M^Y(A^X-A^Y)|\le0.2a(C_*/20+6.9a)\eta|q|^{-1}$. Integrate $|E|\le C_E\eta|q|^{-1}$ against $|u_{r-1}|$. The decaying form is the geometric mean with the separate decay bound $|H^F(x^F,z)|\le Ce^{-\delta|x^F-z|}$ of each chain (Combes--Thomas for the coercive coarse operators \cite{CT73}) and $|x^F-x|\le2\sqrt{13}\eta'$.
\end{proof}

\begin{lemma}[gradients and H\"older bounds at $A=0$]\label{rot:lem:grad}
$|\nabla H^F(x,z)|\le C_\nabla$ for the forward difference of each flow, hence $|H^F(x,z)-H^F(y,z)|\le C|x-y|^{1/2}(e^{-\delta|x-z|/2}+e^{-\delta|y-z|/2})$.
\end{lemma}
\begin{proof}
The difference multiplies the alias-$l$ summand by $d_F(q+l)$ with $|d_F(Q)|\le\min(|Q|,2\sqrt3/\eta')$; the main term is $\le C|q|(1+|q|)^{-2}$ and the aliases give $537a\,\eta$ integrated against $|u_{r-1}|$ with $J_0\le C\log^3(1/\eta)$ (\Cref{rot:lem:CK0}), $\eta\log^3(1/\eta)\le C$. Geometric mean with the decay.
\end{proof}

\begin{lemma}[coarse covariances]\label{rot:lem:Gam}
Let $B^F\defeq(A^F)^{-1}$ and $\Gamma^F\defeq H^FB^FH^{F*}$. \textup{(i)} $C^F=G^F_r+\Gamma^F$ exactly (Woodbury). \textup{(ii)} $\Gamma^F(x,y)=(2\pi)^{-3}\int_{V_1}B^F(p)\hat h^F_x(p)\overline{\hat h^F_y(p)}dp$ with $\hat h^F_x(p)\defeq\sum_zH^F(x,z)e^{-\mathrm ip\cdot z}$; $B^F(p)\le C_B|p|^{-2}$ on $V_1$ and $\sup|\Gamma^F|\le C_1$. \textup{(iii)} $|\Gamma^X(x',y')-\Gamma^Y(x'',y'')|\le C_\Gamma\eta^{1/2}$ for matched pairs. \textup{(iv)} $|\Gamma^F(x,y)-\Gamma^F(x,w)|\le C_2|y-w|^{1/2}$.
\end{lemma}
\begin{proof}
(ii) Parseval on $\Gamma_r$; $A^F(p)^{-1}=a_{r-1}^{-1}+\sum_{l_1}u_{r-1}(p+l_1)^2[\eta^2a^{-1}+S_F(p+l_1)]$ with $S_F(q)\le\frac{\pi^2}{4|q|^2}+31\eta^2$; the term $l_1=0$ gives $\frac{\pi^2}{4|p|^2}+C$ and the others are summable. (iii) $|B^X-B^Y|\le0.198\,\eta^2$ pointwise (\Cref{rot:prop:fc}), $|\hat h^X_{x'}-\hat h^Y_{x''}|\le C'_h\eta^{1/2}$ (\Cref{rot:lem:H}, decaying form); three differences. (iv) \Cref{rot:lem:grad}.
\end{proof}

\begin{lemma}[scale pieces $j\ge2$]\label{rot:lem:slice}
For $j\ge2$, in the unit scale of level $j$ ($\eta_j=5^{1-j}$) and matched points, $|S^X_j(x',y')-S^Y_j(x'',y'')|\le C\eta_j^{1/2}e^{-\delta|x'-y'|}$. In the unit scale of level $r$ ($j<r$): $\le C(5\eta'/\ell_j)^{1/2}\ell_j^{-1}e^{-\delta|x'-y'|/\ell_j}$. The same holds for the photon.
\end{lemma}
\begin{proof}
At level $j$: $|(A^X_j-A^Y_j)(z,w)|\le C\eta_je^{-\delta|z-w|/2}$ (geometric mean of the sup bound $0.20a^2\eta_j^2$ with the decay of each kernel); $C^F_j=(A^F_j+bQ^*Q)^{-1}$ decays for both chains; $C^X_j-C^Y_j=C^X_j(A^Y_j-A^X_j)C^Y_j$. Then $S^X_j-S^Y_j\circ\tau=(H^X_j-H^Y_j\circ\tau)C^X_jH^{X*}_j+H^Y_j\circ\tau(C^X_j-C^Y_j)H^{X*}_j+H^Y_j\circ\tau\,C^Y_j(H^X_j-H^Y_j\circ\tau)^*$; bound the three terms by \Cref{rot:lem:H} at level $j$ and sum the exponentials. Rescale with \cite[(2.20)]{K1}: scalar kernels scale as $\ell^{-1}$ in three dimensions.
\end{proof}

\begin{corollary}[uniform comparison form]\label{rot:cor:unif}
Let $D_j\defeq S^X_j-S^Y_j\circ\tau$ in the unit scale of level $r$ and $\theta\defeq\frac12$. \textup{(a)} Fine slices $0\le j\le r-1$: $|D_j(x',y')|\le C_D\eta'^{\theta}\ell_j^{-1-\theta}e^{-\delta|x'-y'|/\ell_j}$. \textup{(b)} Coarse slices and remainders $j\ge r$: $|D_j|\le C\eta^{1/2}\ell_j^{-1}e^{-\delta|x'-y'|/\ell_j}$ and $|\Gamma^X_n-\Gamma^Y_n\circ\tau|\le C\eta^{1/2}$ at level $n$. \textup{(c)} Given \Cref{rot:prop:BD}(a) for the hybrid flow and \cite[Thm.~3.3]{K1} for the standard flow, the derivative, H\"older and contour versions hold: the contour lines \cite[(3.27)]{K1} along the two paired contour systems differ by the kernel difference along the contour plus the flux of the slice kernel through the sliver between the paired partial contours, a gain $(\eta'/\ell_j)^{1/2}$ on slice $j$; the double contour lines, both ends on contours, have the same relative gain against their size $C\ell_j$.
\end{corollary}
\begin{proof}
(a) For $j\ge2$ this is \Cref{rot:lem:slice}. For $j\in\{0,1\}$ bound $|S^X_j|$ and $|S^Y_j\circ\tau|$ separately by $C\ell_j^{-1}e^{-\delta|\cdot|/\ell_j}$ and use $\ell_j\le5\eta'$, so $1\le5^\theta(\eta'/\ell_j)^\theta$: there is no gain on these lines, and the factor $\eta'^\theta$ is paid for by a degree loss $\theta$ (King's own device at the finest slice, \cite[p.~675]{K1}: ``When $j=0$, we bound each term separately''). (b) Beyond level $r$ both flows apply the same rotated steps; only the input $B^F$ differs, by the relative amount $0.2a\eta^2$ (\Cref{rot:prop:fc}), and every coarse multiplier is built from it by the same formulas with non-negative weights. (c) Interpolation of (a), (b) with the separate bounds of each flow; for the contours, \Cref{rot:lem:pair,rot:lem:flux}: the sliver has area $\le1.5\eta'\ell_{j+1}+33\eta'^2$, against the slice kernel's gradient $\le C\ell_j^{-2}$, plus the end segment $c_\tau\eta'$ and Riemann errors; for $j\in\{0,1\}$ there is no gain and the factor is paid by a degree loss, as in (a).
\end{proof}

\begin{corollary}[the full propagators of the two lattices]\label{rot:cor:C}
On the infinite lattices, for matched pairs and $u=x'-y'$ with $|u|\ge25\eps$ (physical units),
\[
|C^X(x'-y')-C^Y(x''-y'')|\le C_6\,\eps^{1/2}|u|^{-3/2},
\]
and the same for the photon and for $G^F_r$, $\Gamma^F$ separately. For all $u$ the difference is $\le2C_0(|u|+\eps)^{-1}$.
\end{corollary}
\begin{proof}
Pick $n\ge2$ with $\delta_n\le|u|<5\delta_n$ and write $C^F=\sum_{j<n}S^F_j+\Gamma^F_n$. The slices $j\in\{0,1\}$ are each $\le C\eps^{-1}e^{-\delta_*|u|/(5\eps)}\le C'\eps^{1/2}|u|^{-3/2}$; the slices $2\le j<n$ differ by $\le C\delta_j^{-1}(5\eps/\delta_j)^{1/2}e^{-\delta|u|/\delta_j}$ and sum to $\le C\eps^{1/2}|u|^{-3/2}$; the remainder differs by $\le C_\Gamma\delta_n^{-1}(5\eps/\delta_n)^{1/2}$. Pairs stay matched at every level because $\tau$ is defined by the first-step cells.
\end{proof}

Since $C^Y(x''-y'')=C^X(R(x''-y''))$, \Cref{rot:cor:C} is a quantitative rotation estimate for the massive lattice Green's function, with relative error $(\eps/|u|)^{1/2}$, obtained from the two block-spin decompositions alone.
